\chapter{Keeping an eye on things}\label{ch:watch}

\emph{You want to know what all the agents are doing, follow a colleague's work, and not run into
your Claude usage limit by surprise.}

\section{The active counter and the Agents page}\label{sec:agents-page}

\begin{enumerate}
  \item Look at the top bar: the counter \ui{N active} shows how many agents are working right now,
    across all projects you can see. Choose it, or \ui{Agents}, to open the Agents page.
  \item The Agents page (\ui{On the floor}, figure~\ref{fig:agents}) shows the runs in flight as
    cards at the top, and a table of every agent: what it is (chat or pipeline stage), its project
    and issue, its owner, what it is doing at this moment, its model, cost and running time.
  \item Filter with \ui{everyone} or \ui{mine}, and tick \ui{include ended} to see finished agents too.
  \item Choose a run card to open the run page, or an agent to watch its transcript.
\end{enumerate}

\screenshot{18-agents-desktop}{The Agents page: runs in flight and every agent you can see.}{agents}

\section{Watch someone else's session}\label{sec:observe}

Everybody who can read a repository can watch its agents, read-only.

\begin{enumerate}
  \item On a project page, choose another person's agent in the agents \& runs pane, for example
    bob's chat.
  \item The chat pane shows that session live. There is no prompt: only
    the owner can write to their agents (figure~\ref{fig:observing}).
\end{enumerate}

\screenshot{14-observing-desktop}{Observing bob's chat: live, but read-only.}{observing}

\section{Your Claude usage}\label{sec:usage}

Claude limits how much you can use in a five-hour window and in a week. Agern shows how close you
are, for your own account only.

\begin{enumerate}
  \item Look at the usage indicator in the top bar, left of the active counter. It shows the window
    closest to its limit, for example \ui{5h 42\%}. From 80~\% it turns amber and says \ui{high}
    (figure~\ref{fig:usage-high}); at the limit it turns red and says when it resets
    (figure~\ref{fig:usage-limit}).
  \item Choose it to open the \ui{Usage} block in \ui{Settings}, under \ui{Models \& cost}
    (figure~\ref{fig:settings-models} in chapter~\ref{ch:accounts}): every window with its reset time,
    and what you spent through Agern today and this month.
  \item If you are often near the limit, choose cheaper models there (section~\ref{sec:models}).
\end{enumerate}

Claude reports these numbers while one of your agents runs; they are not refreshed in between.
When the limit is reached, your agents cannot go on: a run stops with an error that you can retry
once the limit has reset (section~\ref{sec:retry}), and a release pauses and resumes by itself
(chapter~\ref{ch:releases}).

\screenshot[0.6\linewidth]{46-usage-high-desktop}{The usage indicator at 86~\% of the five-hour window.}{usage-high}

\screenshot[0.6\linewidth]{47-usage-limit-desktop}{The usage indicator when the limit is reached.}{usage-limit}

\whatyouhave
You can see every agent at work, follow your colleagues' sessions, and know how much of your
Claude usage is left. That completes the daily use of Agern.
Upcoming: using Agern from Claude on your own computer is planned and not yet available.
Next, keep Agern itself healthy: chapter~\ref{ch:upgrade}.
